\begin{dang}{Bức xạ cho vân sáng - tối tại vị trí M}
\Anlythuyet{
\begin{itemize}
\item Tại một điểm M bất kì trên màn quan sát, bức xạ \bth{\lambda=\dfrac{a.x_M}{k.D}}\ $(1)$ \textit{\textbf{cho vân sáng}} phải thỏa mãn
\begin{align*}
\boder{0,38\leq \dfrac{a.x_M}{k.D}\leq 0,76} 
\end{align*}
\begin{align*}
\Rightarrow \begin{cases}
\text{Số giá trị k nguyên thoả mãn là số bức xạ cho vân sáng tại M}.\\
\text{Thay k vào \bth{(1)} ta có bước sóng của bức xạ tương ứng cho vân sáng tại M}
\end{cases}
\end{align*}
\item Tại một điểm M bất kì trên màn quan sát, bức xạ \bth{\lambda=\dfrac{a.x_M}{\pbrace{k+\dfrac{1}{2}}.D}}\ $(2)$ \textit{\textbf{cho vân tối}} phải thỏa mãn
\begin{align*}
\boder{0,38\leq \dfrac{a.x_M}{\pbrace{k+\dfrac{1}{2}}.D}\leq 0,76}
\end{align*}
\begin{align*}
\Rightarrow \begin{cases}
\text{Số giá trị k nguyên thoả mãn là số bức xạ cho vân tối tại M}\\
\text{Thay k vào \bth{(2)} ta có bước sóng của bức xạ tương ứng cho vân tối tại M}
\end{cases}
\end{align*}
\end{itemize}
}
\end{dang}